\section{Divide-and-Conquer Experimental Source Code}

This appendix contains the source code used to
implement and experimentally verify the running
time of the divide-and-conquer worst queue
build-up algorithm.

The implementation uses C++14 and was tested
using the MinGW G++ compiler with the
\texttt{-O2} optimization flag.

\subsection{Divide-and-Conquer Implementation}

Listing~\ref{lst:dc-header} contains the C++
implementation of the algorithm, shared by the
demonstration program and the benchmark.

\lstinputlisting[
    language=C++,
    caption={Divide-and-conquer maximum subarray (max\_subarray.h)},
    label={lst:dc-header}
]{code/max_subarray.h}

Listing~\ref{lst:dc-demo} runs the algorithm on
the ten-minute traffic example.

\lstinputlisting[
    language=C++,
    caption={Traffic example demonstration},
    label={lst:dc-demo}
]{code/dc_max_subarray.cpp}

\subsection{Experimental Runtime Measurement}

Listing~\ref{lst:dc-experiment} contains the C++
code used to check correctness, generate random
inputs, measure execution times, and calculate
median running times.

\lstinputlisting[
    language=C++,
    caption={Divide-and-conquer runtime benchmark},
    label={lst:dc-experiment}
]{code/dc_experiment.cpp}

\subsection{Experimental Graph Generation}

Listing~\ref{lst:dc-plot} contains the Python
code used to compare measured execution times
with theoretical $\Theta(n\log n)$ growth.

\lstinputlisting[
    language=Python,
    caption={Python runtime graph generation},
    label={lst:dc-plot}
]{code/dc_plot_results.py}
